% Source: https://artofproblemsolving.com/community/c3h1560683_geometry_problem
% Stable identity: aops-topic-1560683-q01
The parallelogram is spanned by $\vec{AB}=(20,10)$ and $\vec{AD}=(10,y)$, and its area is the absolute value of the determinant:
\[
\left|20\cdot y-10\cdot10\right|=|20y-100| .
\]
Setting $|20y-100|=600$ gives $20y-100=600$ or $20y-100=-600$, so
\[
\boxed{y=35\ \text{or}\ y=-25}.
\]
Both work. In each case $C=B+D-A$, so $C=(30,45)$ or $C=(30,-15)$. The two parallelograms lie on opposite sides of line $AB$, which is why there are two answers.

\begin{asy}
usepackage("amsmath");
size(6.5cm);
defaultpen(fontsize(10pt));
pair A=(0,0),B=(20,10),D1=(10,35),C1=B+D1,D2=(10,-25),C2=B+D2;
fill(A--B--C1--D1--cycle,paleblue);draw(A--B--C1--D1--cycle,blue+linewidth(0.9));
fill(A--B--C2--D2--cycle,palered);draw(A--B--C2--D2--cycle,red+linewidth(0.9));
draw(A--B,linewidth(1.1));
dot(A);dot(B);dot(D1);dot(D2);dot(C1);dot(C2);
label("$A(0,0)$",A,W);label("$B(20,10)$",B,E);
label("$D(10,35)$",D1,NW,blue);label("$C(30,45)$",C1,NE,blue);
label("$D(10,-25)$",D2,SW,red);label("$C(30,-15)$",C2,SE,red);
label("$600$",(A+B+C1+D1)/4,blue);label("$600$",(A+B+C2+D2)/4,red);
label("\mbox{}",(min(currentpicture,user=true)+max(currentpicture,user=true))/2,black); // keeps the text after the figure black when LaTeX typesets the labels inline
\end{asy}

(Check with base and height: $AB=10\sqrt5$, and the distance from $D(10,y)$ to the line $x-2y=0$ is $\dfrac{|10-2y|}{\sqrt5}$; the product is $10|10-2y|=600$, again $y=35$ or $y=-25$.)
